\par\SourceLeaf{165}
Sollen nun aber die Abänderungen für jetzt auf das äusserste Minimum des Zulässigen gestellt werden, so ist die Commission der Meinung, dass diess auf die möglich schonendste Weise durch folgenden Plan geschehen kann, welchen sie daher, bedingungsweise, zur Annahme empfiehlt.

I.\quad Die jährlichen Beiträge werden auf 3 Louisd’or erhöhet.

II.\quad Die Pensionen bleiben für jetzt auf der Höhe von 250 Rthl. bestehen, so lange die Zahl der Witwen\footnote{Es werden hier der Kürze wegen immer nur Witwen genannt, es versteht sich aber von selbst, dass immer eine Waisenpension wie eine Witwenpension gezählt werden muss.}, ohne die Professorin G. mitzuzählen, nicht über 18 hinausgeht, welches die gegenwärtige Anzahl ist. Im entgegengesetzten Falle wird die Pension so bestimmt, dass zu 200 Rthl., als festem Theile, noch die Dividende hinzukommt, welche auf jede einzelne Witwe fällt, indem man 900 Rthl. (als 18 mal 50 Rthl.) unter die vorhandenen gleichmässig vertheilt. Bei der höchsten Witwenzahl, welche bisher vorgekommen ist, nemlich 21 (ohne die Prof. G.) würde also die Pension noch 242 Rthl. 20 Ggr. 7 Pf. Gold betragen.

Man erachtet leicht, dass bei dieser Einrichtung das Deficit noch nicht völlig gedeckt ist: in der That ergibt die Rechnung, dass, den Anfang der Wirksamkeit vom 1. Oct. 1845 an vorausgesetzt, noch 2963 Rthl. ungedeckt bleiben, bei späterm Anfang des erhoheten Beitrags also nach Verhältniss mehr. Wenn jedoch der zur Zeit noch bestehende höhere Zinsgenuss vorerst noch ohne erhebliche Schmälerung fortdauern wird, und sonst keine bedeutenden Verluste eintreten, so könnte man hoffen, dass jener Deficitsrest nach einer mässigen Anzahl von Jahren von selbst zur Deckung gelangen werde. Die Gesellschaft darf jedoch, wenn sie jenen Plan annimmt, sich nicht verhehlen, dass die Witwenkasse unter so sehr knapper Anmessung der Mittel zu den Obliegenheiten einem Schiffe gleicht, welches schwerbelastet in seichtem Fahrwasser geht, und seine Sicherheit nur in fortwährendem wachsamen Sondiren findet.

Die Commission hält daher für unumgänglich nothwendig, dass mit obiger Regulirung noch verbunden werde

III.\quad eine in angemessenen Zeitintervallen (etwa aller 10 oder 5 Jahr) zu wiederholende sorgfältige neue Bilanzrechnung, und dass je nach deren Ergebnissen eventuell weitere Nachhülfe vorbehalten bleibe. Als eine nothwendige vorbereitende Maassregel dazu würde von jetzt an eine solche pünktliche Buchführung über die für die Witwenkasse relevanten personlichen Verhältnisse aller Interessenten, wie bereits an mehrern Stellen der Denkschrift angedeutet ist, einzuführen sein.

Die Commission hält es nicht für nöthig, die Vortheile, welche diese Regulirung darbietet, hier weitläuftig zu entwickeln. und macht nur auf Folgendes aufmerksam. Findet sich nach der ersten oder zweiten Revision, dass die Bilanz sich nicht nur gebessert, sondern viel mehr gebessert habe, als man hatte hoffen können, so wird man die Pensionen weiter erhöhen dürfen von 250 Rthl. auf 260 Rthl. für den Fall einer Witwenzahl unter 18, oder für den Fall der hohen Witwenzahl die ganze Zusatzdividende von 900 Rthl. auf 1080 Rthl. (d. i. 18 mal 60 Rthl.). Findet sich hingegen bei der Revision eine Verschlechterung der Bilanz, so wird man sich einige weitere Anstrengung gefallen lassen aber zugleich um so mehr Glück wünschen müssen, 1846 nicht die Hände in den Schooss gelegt zu haben. In beiden Fällen aber wird man die weitern Maassregeln mit Bewusstsein der Sicherheit oder der Nothwendigkeit treffen können.

Dass mit Annahme dieser Regulirung der Anfang des §. 10 des Regulativs, der in seiner unklaren Fassung als eine Hauptquelle des jetzigen Übels betrachtet werden muss, von selbst wegfällt, braucht nicht erinnert zu werden.

Bei der Anwendung der in diesem Plane enthaltenen Normirung darf abseiten der Witwenkasse kein
